% TOP_PROBLEM: {"id": 3000001, "title": "Achieve global rigidity by pinning nodes", "category_id": 2, "category": {"id": 2, "name": "combinatorics", "display_name": "Combinatorics", "description": "Counting problems, graph theory, discrete structures.", "slug": "combinatorics", "order_index": 2, "created_at": "2026-07-31T15:26:25.671Z"}, "status": "open", "problem_number": "AMR-029-0001", "set_id": 13}
% FIRST_POSTED: 2026-10-01
% ENTRY_KIND: statement_only
% UnsolvedMath ID 3000001; source catalogue code AMR-029-0001
% !TeX program = lualatex
% Definition and sources only; this file is not an LLM solution attempt.
\documentclass[11pt,a4paper]{article}
\usepackage[margin=25mm]{geometry}
\usepackage{fontspec}
\setmainfont{lmroman10-regular.otf}[BoldFont=lmroman10-bold.otf,
 ItalicFont=lmroman10-italic.otf,BoldItalicFont=lmroman10-bolditalic.otf]
\usepackage{amsmath,amssymb,mathtools}
\usepackage{unicode-math}
\setmathfont{latinmodern-math.otf}
\AtBeginDocument{\let\setminus\smallsetminus}
\usepackage{xurl,hyperref}
\hypersetup{colorlinks=true,urlcolor=blue}
\setcounter{secnumdepth}{0}
\setlength{\parindent}{0pt}
\setlength{\parskip}{5pt}
\setlength{\emergencystretch}{3em}
\begin{document}
\section*{Achieve global rigidity by pinning nodes}\label{3000001}
{\small UnsolvedMath ID 3000001 · AMR-029-0001 · Combinatorics · AMR Open Problem Lists}
\subsection{Definitions and mathematical statement}
Let $G=(V,E)$ be a finite simple graph. A placement $p:V\to\mathbb R^2$
is generic when its $2|V|$ coordinates are algebraically independent over
$\mathbb Q$. Two placements are equivalent if their Euclidean edge lengths
agree, and congruent if every pairwise distance agrees. The graph is
generically globally rigid in the plane when every placement equivalent
to a generic placement is congruent to it.

For $S\subseteq V$, let $G+K_S$ be obtained by adding every missing edge
whose two endpoints lie in $S$. Define
\[
 p_2(G)=\min\{|S|:S\subseteq V,\ G+K_S
                     \text{ is generically globally rigid in }\mathbb R^2\}.
\]
The task is to determine this minimum and construct an attaining set,
seeking a polynomial-time algorithm or an effective combinatorial
characterization of the optimum. The input is the graph, without a
special geometric realization.

Adding the clique fixes all relative distances between selected vertices;
it does not fix an absolute coordinate frame. The quantity being minimized
counts selected vertices, not added edges. Selecting all vertices is always
feasible, but provides no optimality guarantee. Global rigidity is essential:
excluding continuous flexes alone does not exclude a second, isolated
noncongruent realization with identical prescribed edge lengths.
\subsection{Short English statement}
How few vertices must be selected so that specifying every distance
among them makes a graph generically globally rigid in the plane?
\subsection{Sources}
\begin{itemize}
\item[S1] ULAM.AI and the UnsolvedMath Contributors, supplied problems.json, record 3000001 (AMR-029-0001), AMR Open Problem Lists. Catalogue identity and source transcription; CC BY 4.0.
\url{https://huggingface.co/datasets/ulamai/UnsolvedMath}.
\item[S2] Egres Open - List of open problems, Open-problem page: Achieve global rigidity by pinning nodes. Original source indexed by AMR-029-0001.
\url{https://oldlemon.cs.elte.hu/egres/open/List_of_open_problems}.
\item[S3] EGRES Open, Achieve global rigidity by pinning nodes. Individual source page and explanatory remarks.
\url{https://lemon.cs.elte.hu/egres/open/Achieve_global_rigidity_by_pinning_nodes}.
\end{itemize}
\end{document}
